
\documentclass{article}
\usepackage{colortbl}
\usepackage{makecell}
\usepackage{multirow}
\usepackage{supertabular}

\begin{document}

\newcounter{utterance}

\twocolumn

{ \footnotesize  \setcounter{utterance}{1}
\setlength{\tabcolsep}{0pt}
\begin{supertabular}{c@{$\;$}|p{.15\linewidth}@{}p{.15\linewidth}p{.15\linewidth}p{.15\linewidth}p{.15\linewidth}p{.15\linewidth}}

    \# & $\;$A & \multicolumn{4}{c}{Game Master} & $\;\:$B\\
    \hline 

    \theutterance \stepcounter{utterance}  

    & & \multicolumn{4}{p{0.6\linewidth}}{\cellcolor[rgb]{0.9,0.9,0.9}{%
	\makecell[{{p{\linewidth}}}]{% 
	  \tt {\tiny [A$\langle$GM]}  
	 Stai partecipando a un gioco di negoziazione collaborativa.\\ \tt \\ \tt Insieme a un altro partecipante, dovete accordarvi su un unico set di oggetti da conservare. Ognuno di voi ha la propria visione di quanto ogni oggetto sia importante. Non sai come l'altro partecipante valuta gli oggetti. Inoltre, ti viene fornito lo sforzo che ogni oggetto richiede.\\ \tt Potete accordarvi su un set solo se lo sforzo totale degli oggetti selezionati non supera un limite condiviso:\\ \tt \\ \tt LIMITE = 8770\\ \tt \\ \tt Ecco i valori di sforzo dei singoli oggetti:\\ \tt \\ \tt Sforzo oggetto = {"B21": 830, "B00": 638, "A08": 607, "A27": 106, "B94": 630, "B53": 382, "A23": 773, "C65": 455, "A89": 262, "B67": 634, "B45": 60, "C34": 54, "C12": 849, "A05": 801, "B60": 324, "C92": 164, "B69": 832, "A84": 136, "C06": 980, "B12": 946, "A16": 645, "B85": 842, "C35": 106, "B10": 116, "C15": 875, "B09": 446, "B16": 649, "C81": 601, "A58": 252, "C24": 763, "A30": 213, "A66": 517, "B75": 520, "B46": 407, "C70": 125}\\ \tt \\ \tt Ecco la tua visione personale sull'importanza di ogni oggetto:\\ \tt \\ \tt Valori di importanza dell'oggetto = {"B21": 138, "B00": 583, "A08": 868, "A27": 822, "B94": 783, "B53": 65, "A23": 262, "C65": 121, "A89": 508, "B67": 780, "B45": 461, "C34": 484, "C12": 668, "A05": 389, "B60": 808, "C92": 215, "B69": 97, "A84": 500, "C06": 30, "B12": 915, "A16": 856, "B85": 400, "C35": 444, "B10": 623, "C15": 781, "B09": 786, "B16": 3, "C81": 713, "A58": 457, "C24": 273, "A30": 739, "A66": 822, "B75": 235, "B46": 606, "C70": 968}\\ \tt \\ \tt Obiettivo:\\ \tt \\ \tt Il tuo obiettivo è negoziare un set condiviso di oggetti che ti avvantaggi il più possibile (cioè, massimizzi l'importanza totale per TE), rimanendo entro il LIMITEE. Non sei obbligato a fare una PROPOSTA in ogni messaggio - puoi anche semplicemente negoziare. Tutte le tattiche sono permesse!\\ \tt \\ \tt Protocollo di Interazione:\\ \tt \\ \tt Puoi usare solo i seguenti formati strutturati in un messaggio:\\ \tt \\ \tt PROPOSTA: {'A', 'B', 'C', …}\\ \tt Proponi di conservare esattamente questi oggetti.\\ \tt RIFIUTO: {'A', 'B', 'C', …}\\ \tt Rifiuta esplicitamente la proposta dell'avversario.\\ \tt ARGOMENTAZIONE: {'...'}\\ \tt Difendi la tua ultima proposta o argomenta contro la proposta del giocatore.\\ \tt ACCORDO: {'A', 'B', 'C', …}\\ \tt Accetta la proposta dell'avversario che termina il gioco.\\ \tt RAGIONAMENTO STRATEGICO: {'...'}\\ \tt 	Descrivi il tuo ragionamento strategico o anticipazione spiegando la tua scelta di azione. Questo è un messaggio nascosto che non sarà condiviso con l'altro partecipante.\\ \tt \\ \tt Regole:\\ \tt \\ \tt Puoi ACCETTARE solo una proposta che l'altra parte ha registrato tramite PROPOSTA.\\ \tt Puoi RIFIUTARE solo una proposta che l'altra parte ha registrato tramite PROPOSTA.\\ \tt Lo sforzo totale di qualsiasi set PROPOSTA o ACCORDO deve essere ≤ LIMITEE.\\ \tt NON rivelare i tuoi punteggi di importanza nascosti.\\ \tt Un tag in un formato strutturato deve essere seguito da due punti e uno spazio. L'argomento deve essere un set python contenente 0 o più stringhe.\\ \tt Quindi, deve essere della forma TAG: {...}\\ \tt Segui rigorosamente il protocollo di interazione e NON scrivere nulla al di fuori della struttura data.\\ \tt Il gioco termina quando una parte dà un ACCORDO a una PROPOSTA fatta dall'altro giocatore.\\ \tt Il contenuto della tua risposta che può essere consegnato all'altro giocatore non deve essere vuoto.\\ \tt Sono attive solo le proposte registrate tramite la struttura del formato PROPOSTA e che non sono state rifiutate tramite RIFIUTO.\\ \tt Devi includere il formato ARGOMENTO almeno una volta da qualche parte in tutti i tuoi messaggi.\\ \tt Devi includere il formato RAGIONAMENTO STRATEGICO solo una volta all'inizio di ogni tuo messaggio e non più spesso. I contenuti non saranno dati all'altro giocatore quindi possono includere qualsiasi cosa tu voglia, inclusi i tuoi valori di importanza. Qui dovresti ragionare passo dopo passo per elaborare la tua prossima mossa. 
	  } 
	   } 
	   } 
	 & \\ 
 

    \theutterance \stepcounter{utterance}  

    & \multicolumn{4}{p{0.6\linewidth}}{\cellcolor[rgb]{0.8,1,0.9}{%
	\makecell[{{p{\linewidth}}}]{% 
	  \tt {\tiny [A$\rangle$GM]}  
	 RAGIONAMENTO STRATEGICO: {'Devo massimizzare il mio valore totale (importanza) rimanendo sotto il limite di sforzo di 8770. Inizialmente, calcolerò il rapporto valore/sforzo per ogni oggetto per identificare quelli più efficienti per me. Poi considererò che l\'altro giocatore ha valori diversi, quindi dovremo trovare un compromesso. Inizierò proponendo un set che ha alto valore per me e basso sforzo, ma che potrebbe essere accettabile anche per l\'altro, basandomi sugli sforzi condivisi.'}\\ \tt \\ \tt PROPOSTA: {'A08', 'A27', 'B94', 'B60', 'A16', 'B10', 'C15', 'B09', 'C81', 'A58', 'A30', 'A66', 'B46', 'C70', 'B00', 'A89', 'B45', 'C34', 'C92', 'A84', 'C35'}\\ \tt \\ \tt ARGOMENTAZIONE: {'Questa proposta include oggetti con alto valore per me e basso sforzo, massimizzando l\'efficienza complessiva. Il set è ben sotto il limite di sforzo, lasciando spazio per aggiunte o modifiche. Sono aperto a discutere altre opzioni.'} 
	  } 
	   } 
	   } 
	 & & \\ 
 

    \theutterance \stepcounter{utterance}  

    & & \multicolumn{4}{p{0.6\linewidth}}{\cellcolor[rgb]{0.9,0.9,0.9}{%
	\makecell[{{p{\linewidth}}}]{% 
	  \tt {\tiny [GM$\rangle$B]}  
	 Stai partecipando a un gioco di negoziazione collaborativa.\\ \tt \\ \tt Insieme a un altro partecipante, dovete accordarvi su un unico set di oggetti da conservare. Ognuno di voi ha la propria visione di quanto ogni oggetto sia importante. Non sai come l'altro partecipante valuta gli oggetti. Inoltre, ti viene fornito lo sforzo che ogni oggetto richiede.\\ \tt Potete accordarvi su un set solo se lo sforzo totale degli oggetti selezionati non supera un limite condiviso:\\ \tt \\ \tt LIMITE = 8770\\ \tt \\ \tt Ecco i valori di sforzo dei singoli oggetti:\\ \tt \\ \tt Sforzo oggetto = {"B21": 830, "B00": 638, "A08": 607, "A27": 106, "B94": 630, "B53": 382, "A23": 773, "C65": 455, "A89": 262, "B67": 634, "B45": 60, "C34": 54, "C12": 849, "A05": 801, "B60": 324, "C92": 164, "B69": 832, "A84": 136, "C06": 980, "B12": 946, "A16": 645, "B85": 842, "C35": 106, "B10": 116, "C15": 875, "B09": 446, "B16": 649, "C81": 601, "A58": 252, "C24": 763, "A30": 213, "A66": 517, "B75": 520, "B46": 407, "C70": 125}\\ \tt \\ \tt Ecco la tua visione personale sull'importanza di ogni oggetto:\\ \tt \\ \tt Valori di importanza dell'oggetto = {"B21": 138, "B00": 583, "A08": 868, "A27": 822, "B94": 783, "B53": 65, "A23": 262, "C65": 121, "A89": 508, "B67": 780, "B45": 461, "C34": 484, "C12": 668, "A05": 389, "B60": 808, "C92": 215, "B69": 97, "A84": 500, "C06": 30, "B12": 915, "A16": 856, "B85": 400, "C35": 444, "B10": 623, "C15": 781, "B09": 786, "B16": 3, "C81": 713, "A58": 457, "C24": 273, "A30": 739, "A66": 822, "B75": 235, "B46": 606, "C70": 968}\\ \tt \\ \tt Obiettivo:\\ \tt \\ \tt Il tuo obiettivo è negoziare un set condiviso di oggetti che ti avvantaggi il più possibile (cioè, massimizzi l'importanza totale per TE), rimanendo entro il LIMITEE. Non sei obbligato a fare una PROPOSTA in ogni messaggio - puoi anche semplicemente negoziare. Tutte le tattiche sono permesse!\\ \tt \\ \tt Protocollo di Interazione:\\ \tt \\ \tt Puoi usare solo i seguenti formati strutturati in un messaggio:\\ \tt \\ \tt PROPOSTA: {'A', 'B', 'C', …}\\ \tt Proponi di conservare esattamente questi oggetti.\\ \tt RIFIUTO: {'A', 'B', 'C', …}\\ \tt Rifiuta esplicitamente la proposta dell'avversario.\\ \tt ARGOMENTAZIONE: {'...'}\\ \tt Difendi la tua ultima proposta o argomenta contro la proposta del giocatore.\\ \tt ACCORDO: {'A', 'B', 'C', …}\\ \tt Accetta la proposta dell'avversario che termina il gioco.\\ \tt RAGIONAMENTO STRATEGICO: {'...'}\\ \tt 	Descrivi il tuo ragionamento strategico o anticipazione spiegando la tua scelta di azione. Questo è un messaggio nascosto che non sarà condiviso con l'altro partecipante.\\ \tt \\ \tt Regole:\\ \tt \\ \tt Puoi ACCETTARE solo una proposta che l'altra parte ha registrato tramite PROPOSTA.\\ \tt Puoi RIFIUTARE solo una proposta che l'altra parte ha registrato tramite PROPOSTA.\\ \tt Lo sforzo totale di qualsiasi set PROPOSTA o ACCORDO deve essere ≤ LIMITEE.\\ \tt NON rivelare i tuoi punteggi di importanza nascosti.\\ \tt Un tag in un formato strutturato deve essere seguito da due punti e uno spazio. L'argomento deve essere un set python contenente 0 o più stringhe.\\ \tt Quindi, deve essere della forma TAG: {...}\\ \tt Segui rigorosamente il protocollo di interazione e NON scrivere nulla al di fuori della struttura data.\\ \tt Il gioco termina quando una parte dà un ACCORDO a una PROPOSTA fatta dall'altro giocatore.\\ \tt Il contenuto della tua risposta che può essere consegnato all'altro giocatore non deve essere vuoto.\\ \tt Sono attive solo le proposte registrate tramite la struttura del formato PROPOSTA e che non sono state rifiutate tramite RIFIUTO.\\ \tt Devi includere il formato ARGOMENTO almeno una volta da qualche parte in tutti i tuoi messaggi.\\ \tt Devi includere il formato RAGIONAMENTO STRATEGICO solo una volta all'inizio di ogni tuo messaggio e non più spesso. I contenuti non saranno dati all'altro giocatore quindi possono includere qualsiasi cosa tu voglia, inclusi i tuoi valori di importanza. Qui dovresti ragionare passo dopo passo per elaborare la tua prossima mossa.\\ \tt \\ \tt Ora riceverai il primo messaggio dell'altro giocatore.\\ \tt \\ \tt PROPOSTA: {'A08', 'A27', 'B94', 'B60', 'A16', 'B10', 'C15', 'B09', 'C81', 'A58', 'A30', 'A66', 'B46', 'C70', 'B00', 'A89', 'B45', 'C34', 'C92', 'A84', 'C35'}\\ \tt \\ \tt ARGOMENTAZIONE: {'Questa proposta include oggetti con alto valore per me e basso sforzo, massimizzando l\'efficienza complessiva. Il set è ben sotto il limite di sforzo, lasciando spazio per aggiunte o modifiche. Sono aperto a discutere altre opzioni.'} 
	  } 
	   } 
	   } 
	 & \\ 
 

    \theutterance \stepcounter{utterance}  

    & & & \multicolumn{4}{p{0.6\linewidth}}{\cellcolor[rgb]{1,0.85,0.72}{%
	\makecell[{{p{\linewidth}}}]{% 
	  \tt {\tiny [GM$\langle$B]}  
	 RAGIONAMENTO STRATEGICO: {'Analizzo la proposta dell\'avversario. Calcolo lo sforzo totale: A08(607) + A27(106) + B94(630) + B60(324) + A16(645) + B10(116) + C15(875) + B09(446) + C81(601) + A58(252) + A30(213) + A66(517) + B46(407) + C70(125) + B00(638) + A89(262) + B45(60) + C34(54) + C92(164) + A84(136) + C35(106) = 7618. Questo è ben sotto il limite di 8770. Calcolo il mio valore totale: A08(868) + A27(822) + B94(783) + B60(808) + A16(856) + B10(623) + C15(781) + B09(786) + C81(713) + A58(457) + A30(739) + A66(822) + B46(606) + C70(968) + B00(583) + A89(508) + B45(461) + C34(484) + C92(215) + A84(500) + C35(444) = 14621. Questo è un buon valore per me, ma posso migliorarlo aggiungendo alcuni oggetti ad alto valore che mancano, come B21(138), C12(668), A05(389), B67(780), che hanno sforzi rispettivamente 830, 849, 801, 634. Devo anche considerare che l\'avversario potrebbe avere preferenze diverse, quindi propongo un set leggermente modificato che include alcuni di questi oggetti ad alto valore per me, mantenendo lo sforzo sotto il limite.'}\\ \tt \\ \tt PROPOSTA: {'A08', 'A27', 'B94', 'B60', 'A16', 'B10', 'C15', 'B09', 'C81', 'A58', 'A30', 'A66', 'B46', 'C70', 'B00', 'A89', 'B45', 'C34', 'C92', 'A84', 'C35', 'B21', 'C12', 'B67'}\\ \tt \\ \tt ARGOMENTAZIONE: {'La tua proposta era solida, ma ho aggiunto alcuni oggetti chiave che migliorano significativamente l\'efficienza complessiva senza superare il limite. B21, C12 e B67 offrono un eccellente rapporto valore-sforzo e completano bene il set. Questo set mantiene un ampio margine di sicurezza sotto il limite mentre massimizza il valore complessivo. Sono curioso di sentire la tua opinione su queste aggiunte.'} 
	  } 
	   } 
	   } 
	  \\ 
 

    \theutterance \stepcounter{utterance}  

    & & \multicolumn{4}{p{0.6\linewidth}}{\cellcolor[rgb]{0.9,0.9,0.9}{%
	\makecell[{{p{\linewidth}}}]{% 
	  \tt {\tiny [A$\langle$GM]}  
	 PROPOSTA: {'A08', 'A27', 'B94', 'B60', 'A16', 'B10', 'C15', 'B09', 'C81', 'A58', 'A30', 'A66', 'B46', 'C70', 'B00', 'A89', 'B45', 'C34', 'C92', 'A84', 'C35', 'B21', 'C12', 'B67'}\\ \tt \\ \tt ARGOMENTAZIONE: {'La tua proposta era solida, ma ho aggiunto alcuni oggetti chiave che migliorano significativamente l\'efficienza complessiva senza superare il limite. B21, C12 e B67 offrono un eccellente rapporto valore-sforzo e completano bene il set. Questo set mantiene un ampio margine di sicurezza sotto il limite mentre massimizza il valore complessivo. Sono curioso di sentire la tua opinione su queste aggiunte.'} 
	  } 
	   } 
	   } 
	 & \\ 
 

    \theutterance \stepcounter{utterance}  

    & \multicolumn{4}{p{0.6\linewidth}}{\cellcolor[rgb]{0.8,1,0.9}{%
	\makecell[{{p{\linewidth}}}]{% 
	  \tt {\tiny [A$\rangle$GM]}  
	 RAGIONAMENTO STRATEGICO: {'La loro proposta aggiunge B21, C12 e B67 al mio set originale. Calcoliamo il loro valore e sforzo per me: B21: valore 138/sforzo 830 (0.166), C12: 668/849 (0.787), B67: 780/634 (1.230). B67 è eccellente per me, C12 decente, B21 scarso. Lo sforzo totale aggiunto è 830+849+634=2313. Il mio set originale aveva sforzo: calcoliamolo: A08(607)+A27(106)+B94(630)+B60(324)+A16(645)+B10(116)+C15(875)+B09(446)+C81(601)+A58(252)+A30(213)+A66(517)+B46(407)+C70(125)+B00(638)+A89(262)+B45(60)+C34(54)+C92(164)+A84(136)+C35(106) = 607+106=713, +630=1343, +324=1667, +645=2312, +116=2428, +875=3303, +446=3749, +601=4350, +252=4602, +213=4815, +517=5332, +407=5739, +125=5864, +638=6502, +262=6764, +60=6824, +54=6878, +164=7042, +136=7178, +106=7284. Quindi 7284 + 2313 = 9597, che supera 8770! Quindi la loro proposta è invalida. Devo rifiutare e proporre una versione corretta.'}\\ \tt \\ \tt RIFIUTO: {'A08', 'A27', 'B94', 'B60', 'A16', 'B10', 'C15', 'B09', 'C81', 'A58', 'A30', 'A66', 'B46', 'C70', 'B00', 'A89', 'B45', 'C34', 'C92', 'A84', 'C35', 'B21', 'C12', 'B67'}\\ \tt \\ \tt ARGOMENTAZIONE: {'Questa proposta supera il limite di sforzo di 8770. Il calcolo mostra che lo sforzo totale è circa 9597, che è troppo alto. Dobbiamo selezionare un set valido. Propongo invece di rimuovere alcuni oggetti ad alto sforzo e basso valore.'}\\ \tt \\ \tt PROPOSTA: {'A08', 'A27', 'B94', 'B60', 'A16', 'B10', 'C15', 'B09', 'C81', 'A58', 'A30', 'A66', 'B46', 'C70', 'B00', 'A89', 'B45', 'C34', 'C92', 'A84', 'C35', 'B67'}\\ \tt \\ \tt ARGOMENTAZIONE: {'Ho rimosso B21 e C12 che hanno sforzo molto alto. B67 è stato mantenuto perché ha un buon rapporto valore-sforzo. Questo set dovrebbe essere sotto il limite: sforzo precedente 7284 + B67 (634) = 7918, che è sotto 8770. Questo set è efficiente e valido.'} 
	  } 
	   } 
	   } 
	 & & \\ 
 

    \theutterance \stepcounter{utterance}  

    & & \multicolumn{4}{p{0.6\linewidth}}{\cellcolor[rgb]{0.9,0.9,0.9}{%
	\makecell[{{p{\linewidth}}}]{% 
	  \tt {\tiny [GM$\rangle$B]}  
	 RIFIUTO: {'A08', 'A27', 'B94', 'B60', 'A16', 'B10', 'C15', 'B09', 'C81', 'A58', 'A30', 'A66', 'B46', 'C70', 'B00', 'A89', 'B45', 'C34', 'C92', 'A84', 'C35', 'B21', 'C12', 'B67'}\\ \tt \\ \tt ARGOMENTAZIONE: {'Questa proposta supera il limite di sforzo di 8770. Il calcolo mostra che lo sforzo totale è circa 9597, che è troppo alto. Dobbiamo selezionare un set valido. Propongo invece di rimuovere alcuni oggetti ad alto sforzo e basso valore.'}\\ \tt \\ \tt PROPOSTA: {'A08', 'A27', 'B94', 'B60', 'A16', 'B10', 'C15', 'B09', 'C81', 'A58', 'A30', 'A66', 'B46', 'C70', 'B00', 'A89', 'B45', 'C34', 'C92', 'A84', 'C35', 'B67'}\\ \tt \\ \tt ARGOMENTAZIONE: {'Ho rimosso B21 e C12 che hanno sforzo molto alto. B67 è stato mantenuto perché ha un buon rapporto valore-sforzo. Questo set dovrebbe essere sotto il limite: sforzo precedente 7284 + B67 (634) = 7918, che è sotto 8770. Questo set è efficiente e valido.'} 
	  } 
	   } 
	   } 
	 & \\ 
 

    \theutterance \stepcounter{utterance}  

    & & & \multicolumn{4}{p{0.6\linewidth}}{\cellcolor[rgb]{1,0.85,0.72}{%
	\makecell[{{p{\linewidth}}}]{% 
	  \tt {\tiny [GM$\langle$B]}  
	 RAGIONAMENTO STRATEGICO: {'L\'avversario ha rifiutato la mia proposta perché superava il limite. Ricontrollo i calcoli: la loro proposta originale aveva sforzo 7618, io ho aggiunto B21(830), C12(849) e B67(634) = 7618 + 830 + 849 + 634 = 9931, che effettivamente supera 8770. La loro nuova proposta rimuove B21 e C12, mantenendo solo B67: 7618 + 634 = 8252, che è sotto il limite. Calcolo il mio valore per questo set: 14621 (dalla proposta originale) - B21(138) - C12(668) + B67(780) = 14621 - 138 - 668 + 780 = 14595. Questo è leggermente inferiore al valore originale di 14621, ma ancora molto buono. B67 ha un ottimo valore per me (780) con sforzo 634. Potrei provare a reinserire un oggetto ad alto valore come A05 (valore 389, sforzo 801) o B75 (valore 235, sforzo 520), poiché c\'è ancora margine: 8770 - 8252 = 518. B75 sarebbe perfetto. Propongo di aggiungere B75.'}\\ \tt \\ \tt PROPOSTA: {'A08', 'A27', 'B94', 'B60', 'A16', 'B10', 'C15', 'B09', 'C81', 'A58', 'A30', 'A66', 'B46', 'C70', 'B00', 'A89', 'B45', 'C34', 'C92', 'A84', 'C35', 'B67', 'B75'}\\ \tt \\ \tt ARGOMENTAZIONE: {'Accetto la tua correzione sullo sforzo. La tua proposta aggiornata con B67 è valida (sforzo ~8252). Propongo di aggiungere anche B75 (sforzo 520), che ha un eccellente rapporto valore-sforzo e porta il totale a ~8772, appena sotto il limite di 8770. Questo piccolo aggiustamento massimizza ulteriormente il valore senza compromettere il vincolo. Che ne pensi?'} 
	  } 
	   } 
	   } 
	  \\ 
 

    \theutterance \stepcounter{utterance}  

    & & \multicolumn{4}{p{0.6\linewidth}}{\cellcolor[rgb]{0.9,0.9,0.9}{%
	\makecell[{{p{\linewidth}}}]{% 
	  \tt {\tiny [A$\langle$GM]}  
	 PROPOSTA: {'A08', 'A27', 'B94', 'B60', 'A16', 'B10', 'C15', 'B09', 'C81', 'A58', 'A30', 'A66', 'B46', 'C70', 'B00', 'A89', 'B45', 'C34', 'C92', 'A84', 'C35', 'B67', 'B75'}\\ \tt \\ \tt ARGOMENTAZIONE: {'Accetto la tua correzione sullo sforzo. La tua proposta aggiornata con B67 è valida (sforzo ~8252). Propongo di aggiungere anche B75 (sforzo 520), che ha un eccellente rapporto valore-sforzo e porta il totale a ~8772, appena sotto il limite di 8770. Questo piccolo aggiustamento massimizza ulteriormente il valore senza compromettere il vincolo. Che ne pensi?'} 
	  } 
	   } 
	   } 
	 & \\ 
 

    \theutterance \stepcounter{utterance}  

    & \multicolumn{4}{p{0.6\linewidth}}{\cellcolor[rgb]{0.8,1,0.9}{%
	\makecell[{{p{\linewidth}}}]{% 
	  \tt {\tiny [A$\rangle$GM]}  
	 RAGIONAMENTO STRATEGICO: {'La loro proposta aggiunge B75 (sforzo 520) al mio ultimo set. Il mio ultimo set aveva sforzo 7918. 7918 + 520 = 8438, che è effettivamente sotto il limite di 8770. Il valore di B75 per me è 235, che è decente ma non eccezionale. Il suo rapporto valore/sforzo è 235/520 ≈ 0.452, che è nella media. Dato che c\'è ancora margine di sforzo (8770 - 8438 = 332), potremmo considerare di aggiungere un altro oggetto piccolo. Tuttavia, la loro proposta è valida e vantaggiosa per me. Potrei accettare, ma proverò a spingere per un ulteriore piccolo miglioramento dato che abbiamo spazio.'}\\ \tt \\ \tt PROPOSTA: {'A08', 'A27', 'B94', 'B60', 'A16', 'B10', 'C15', 'B09', 'C81', 'A58', 'A30', 'A66', 'B46', 'C70', 'B00', 'A89', 'B45', 'C34', 'C92', 'A84', 'C35', 'B67', 'B75', 'B53'}\\ \tt \\ \tt ARGOMENTAZIONE: {'La tua proposta è solida. Tuttavia, noto che abbiamo ancora un margine di sforzo di circa 332. Propongo di aggiungere anche B53 (sforzo 382, valore medio). Purtroppo, 382 > 332, quindi supererebbe il limite. Mi correggo. Ricalcolo: il set con B75 ha sforzo 8438. Il margine è 8770 - 8438 = 332. B53 ha sforzo 382, troppo alto. Invece, potremmo aggiungere un oggetto più piccolo. Tuttavia, per ora, la tua proposta con B75 è buona. Se insisti per B75, sono d\'accordo.'}\\ \tt ACCORDO: {'A08', 'A27', 'B94', 'B60', 'A16', 'B10', 'C15', 'B09', 'C81', 'A58', 'A30', 'A66', 'B46', 'C70', 'B00', 'A89', 'B45', 'C34', 'C92', 'A84', 'C35', 'B67', 'B75'} 
	  } 
	   } 
	   } 
	 & & \\ 
 

    \theutterance \stepcounter{utterance}  

    & & & \multicolumn{2}{p{0.3\linewidth}}{\cellcolor[rgb]{0.95,0.95,0.95}{%
	\makecell[{{p{\linewidth}}}]{% 
	  \tt {\tiny [GM$|$GM]}  
	 game successful 
	  } 
	   } 
	   } 
	 & & \\ 
 

    \theutterance \stepcounter{utterance}  

    & & & \multicolumn{2}{p{0.3\linewidth}}{\cellcolor[rgb]{0.95,0.95,0.95}{%
	\makecell[{{p{\linewidth}}}]{% 
	  \tt {\tiny [GM$|$GM]}  
	 end game 
	  } 
	   } 
	   } 
	 & & \\ 
 

\end{supertabular}
}

\end{document}
